\section{Conclusion}

We presented TSSMRBench, a benchmark for temporal semantic state memory retrieval. Its target is precise: given an evolving memory trajectory, can a memory system retrieve the version memory that matches the queried temporal semantic state, rather than merely a relevant topic or the most recent state.

To study this ability, we built a benchmark with 300 release-evolution scenarios, 9{,}000 temporal semantic states, and 900 questions spanning single-state lookup, cross-version comparison, and temporal ordering. This design turns stage-correct version retrieval into a directly testable benchmark problem and separates it from ordinary topical retrieval.

Our experiments show that current memory systems remain far stronger on single-state lookup than on questions that require multiple temporally related states. FAISS is the strongest automatic baseline because it retrieves complete version-level memory units effectively, but it still degrades when several nearby versions are all relevant. Mem0 is limited by fact-level extraction together with memory merging and update, which often weakens version-state fidelity. Graphiti benefits from structured retrieval on some harder questions, yet it still frequently retrieves the right neighborhood without recovering the full stage-correct version set. These results show that retrieving evolving memories is not enough: the central challenge is preserving and recovering the correct version state.

The main design implication is clear. Future memory systems should not only store evolving information, but also preserve version history, return coherent state-level evidence, distinguish among semantically similar stages, and do so efficiently in large shared memory pools. TSSMRBench provides a benchmark for measuring that capability directly.
